% Last Updated: March 28, 2023
	\documentclass[a4paper, twoside]{article}
	% Margins
	\usepackage{geometry}
		\geometry{
			% showframe,
			left   = 20mm,
			right  = 20mm,
			top    = 20mm,
			bottom = 20mm
		}

	% Headers and footers
	\usepackage{fancyhdr}
		\pagestyle{fancy}
		\fancyhf{}
		\fancyhead[LE, RO]{\thepage}
		\fancyhead[CE]{\textsc{Cong Wen}}
		\fancyhead[CO]{\textsc\rightmark}

		\setlength{\headheight}{15pt}
		\renewcommand\headrulewidth{0pt}








\input{../universal-pre}

\title{\tbf{Lawrence-Venkatesh method}}
\author{Cong Wen, 09/2023}
\date{}
% Last Updated: March 28, 2023



\begin{document}
\maketitle
\thispagestyle{empty}
\begin{abstract}
	This is my self-use notes\footnote{Last updated on June 02, 2023} on \tbf{Lawrence-Venkatesh method}.


	Please reach out if you find anything incorrect.
\end{abstract}
\tableofcontents
% ----------------------------------------------------------------------------------------------------

\sct{Lawrence-Venkatesh}

Based on their proof \cite{LV2020}. See also Lawrence's NSF grant \href{https://www.nsf.gov/awardsearch/showAward?AWD_ID=2313466}{description}.


\begin{ups}
	Lawrence-Venkatesh reproved Mordell conjecture. According to \cite[Section~1.5]{LV2020}, comparing their proof to the proof of Faltings, their proof is (1)stronger in that it is easier to apply to higher dimensional varieties and (2)weaker in that it can say nothing about the Tate and Shafarevich conjectures. Regarding the proof of Mordell conjecture, both methods made crucial use of the Kodaira-Parshin family.


	Let $K$ be a number field, $Y$ a smooth projective curve over $K$ of genus $g\gqs2$, take a smooth projective morphism $f:X\ar Y$, then for good primes $p$ and each $y$, there is $\rho_{y}:G_{K}\ar\Aut(\mathrm{H}_{\text{{\'e}t}}^{\blt}((X_{y})_{\oL{K}},\bQ_{p}))$. Then Mordell conjecture can be reduced to the following for a suitable choice of $f$ and finite place $v|p$.
	\begin{enr}
		\item $\rho_{x}$ is semisimple for almost all $x\in X(K)$,
		\item The following map is finite-to-one,
		\begin{align*}
		  Y(K)&\ar \{\text{$G_{K_{v}}$-representations}\}/\cong \\
		  y&\amt \rho_{y}|_{G_{K_{v}}}
		\end{align*}
	\end{enr}

	To deal with (1), (Note that for abelian varieties as in Faltings' proof, de Rham cohomology in all degrees are made up by those in degree $1$, which are essentially just Tate modules, so the rational Tate module $V_{\ell}(A)$ is a semisimple $G_{K}$-modules really means a lot in this setting.) one uses $p$-adic Hodge theory to put constraints on the Hodge weights of a global representation, such purity argument is reminiscent of how Faltings used Raynaud's result (which is actually the key to Faltings himself).

	To deal with (2), (which is purely local in nature), note $\rho_{x}|_{G_{K_{v}}}$ is crystalline, $p$-adic Hodge theory guarantees crystalline representations fully faithfully embed in filtered $\phi$-modules over $K_{v}$, then it suffices to show
	\begin{align*}
	  Y(K_{v})&\ar \{\text{filtered $\phi$-modules over $K_{v}$}\}/\cong \\
	  y&\amt (\mathrm{H}_{\text{dR}}^{\blt}(Y_{x}/K_{v}),\Fil\xB,\vfi)
	\end{align*}
	is finite-to-one. The Gauss-Manin connection identifies cohomology of nearby fibers, which respects $\vfi$, but does not respect $\Fil\xB$, so choose a small $p$-adic disk $\Oga\sbs X(K_{v})$, there is a $K_{v}$-analytic period map $\Phi_{v}:\Oga\ar\cF(K_{v})$, where $\cF$ is a suitable flag variety parametrising filtrations of $H:=\mathrm{H}_{\text{dR}}^{\blt}(Y_{x_{0}}/K_{v})$. So the following are equivalent
	\begin{itm}
		\item $y_{1},y_{2}$ are in the same fiber.
		\item $(H,\Phi_{v}(x_{1}),\vfi), (H,\Phi_{v}(x_{2}),\vfi)$ are isomorphic as filtered $\phi$-modules.
		\item $\Phi_{v}(x_{1})$ can be transformed to $\Phi_{v}(x_{2})$ via an element in $Z(\vfi)$.
	\end{itm}
	Using this, to prove the above is finite-to-one boils down to prove $\Phi_{v}\xI(\Phi_{v}(\Oga)\cap Z(\vfi)\Phi_{v}(y_{1}))$ is finite. This in turn reduces to show
	\begin{enr}
		\item the $Z(\vfi)$-orbit is a proper subvariety.
		\item the $\Phi_{v}$-image is Zariski dense.
	\end{enr}
	Now choose $f:X\ar Y$ to be the Kodaira-Parshin family. To show $Z(\vfi)$-orbit is small, one can just enlarge the base field; to show $\Phi_{v}$-image is large, one can pass to the complex period map and do explicit topological computation of the monodromy action of $\pi_{1}(X)$. Dealing with these two problems takes up large parts of their paper.


	The key statements of Lawrence-Venkatesh are collected below for convenience.
\end{ups}


% \begin{stp}
% 	Let $K$ be a number field, $S$ a finite set of finite places of $K$. Let $\pi:X\ar Y$ be a proper smooth morphism over $K$ extending to a proper smooth morphism $\pi:\cX\ar\cY$ over $\cO_{K,S}$. Fix an infinite place $\ita:K\ahr\bC$, and a finite place $v:K\ahr K_{v}$ which is unramified above an odd prime $p$ and all conjugates of $v$ do not lie in $S$. By enlarging $S$ when necessary, the relative algebraic de Rham cohomology $\cH^{q}=R^{q}\pi\zS\Oga_{\cX/\cY}\xB$ are sheaves of locally free $\cO_{\cY}$-modules whose generic fiber over $K$ is equipped with the Gauss-Manin connection which can extend to $\nbl:\cH^{q}\ar\cH^{q}\ot\Oga_{\cY/\cO_{K,S}}^{1}$ over $\cO_{K,S}$.


% 	Fix $y_{0}\in\cY(\cO_{K,S})$. Then $V=\mathrm{H}_{\text{dR}}^{q}(\cX_{y_{0}}/K)$ is a finite dimensional $K$-vector space, let $\cF$ be the flag variety over $K$ parametrizing flags in $V$ with the same dimensional data as $\Fil\xB V$. Locally, coefficients of a flat section with respect to Gauss-Manin connection of $\cH^{q}$ can be written as power series in $K\dS{z_{1},\ldots,z_{m}}$, evaluating the flat sections at $y\in\cY(\cO_{K,S})$ which is close to $y_{0}$ gives isomorphism of cohomology of nearby fibers, hence period maps, in the following sense:
% 	\begin{itm}
% 		\item Let $\Oga_{v}=\{y\in\cY(\cO_{K,S})\mid \oL{y}=\oL{y_{0}}\in\cY(\kpa_{v})\}$, and $\oL{\cX}=\Img(\cX_{y_{0}}(\cO_{K,S})\ar\cX_{y_{0}}(K_{v}))^{\text{Zar}}$. Then Gauss-Manin connection induces an isomorphism respecting the crystalline comparison
% 		\[
% 			\begin{tikzcd}
% 				\mathrm{H}_{\text{dR}}^{q}(\cX_{y}/K_{v})\ar[rr, "\cong"]\ar[rd, "\cong"] && \mathrm{H}_{\text{dR}}^{q}(\cX_{y_{0}}/K_{v})\ar[ld, "\cong"']\\
% 				& \mathrm{H}_{\text{cris}}^{q}(\oL{\cX})\ot K_{v}&
% 			\end{tikzcd}
% 		\]
% 		and a $K_{v}$-analytic period map
% 		\[
% 			\Phi_{v}:\Oga_{v}\ar\cF(K_{v}).
% 		\]

% 		\item Let $\Oga_{\bC}$ be a contractible analytic neighborhood of $y_{0}\in Y_{\bC}^{\text{an}}$. Then Gauss-Manin connection induces
% 		\[
% 			\mathrm{H}_{\text{dR}}^{q}((\cX_{y_{1}})_{\bC}/\bC)\cong \mathrm{H}_{\text{dR}}^{q}((\cX_{y_{2}})_{\bC}/\bC),
% 		\]
% 		and a complex period map equivariant for the monodromy action of $\pi_{1}(Y_{\bC}^{\text{an}},y_{0})\ar\Aut(V_{\bC})$,
% 		\[
% 			\Phi_{\bC}:\Oga_{\bC}\ar\cF(\bC).
% 		\]
% 	\end{itm}

% \end{stp}


% \begin{thm}[{\cite[Lemma~3.1, Lemma~3.2, Lemma~3.3]{LV2020}}]
% 	Let $\Gma=\Img(\pi_{1}(Y_{\bC}^{\text{an}},y_{0})\ar\Aut\cF(\bC))^{\text{Zar}}$. Then
% 	\[
% 		\Dim_{\bC}(\Gma\cdot\Phi_{\bC}(y_{0}))\lqs \Dim_{\bC}\Phi_{\bC}(\Oga_{\bC})=\Dim_{K_{v}}\Phi_{v}(\Oga_{v}).
% 	\]
% \end{thm}

% \begin{thm}[{\cite[Proposition~3.4]{LV2020}}]
% 	Suppose $\Dim_{K_{v}}(Z(\phi_{v}^{[K_{v}:\bQ_{p}]}))<\Dim_{\bC}(\Gma\cdot\Phi_{\bC}(y_{0}))$, then
% 	\[
% 		\{y\in\Oga_{v}\mid \text{$\rho_{y}$ is semisimple}\}\sbs\Oga_{v}
% 	\]
% 	is contained in a proper $K_{v}$-analytic subvariety of $\Oga_{v}$.
% \end{thm}

% \prg{S-unit equation, 4.2, 4.3, 4.4}

% \begin{thm}[{\cite[Theorem~4.1]{LV2020}}]
% 	Notations as above, one has
% 	\[
% 		\#\bmat{t\in\cO_{K,S}\xT\mid 1-t\in\cO_{K,S}\xT}<+\ift.
% 	\]
% \end{thm}

% \begin{prf}[Sketch of the proof]
% 	First, one can assume $K\sps\mu_{8}$, $S$ contains all primes above $2$. The statement can be reduced to the following case:

% 	(\cite[Lemma~4.2]{LV2020}) Fix a cyclic extension $L/K$, a place $v\nin S$ such that
% 	\begin{itm}
% 		\item $\Gal(L/K)=\sA{\Frob_{v}}$,
% 		\item the prime $p$ below $v$ is unramified in $K$,
% 		\item no prime of $S$ lies above $p$,
% 	\end{itm}
% 	let $m$ be the largest power of $2$ dividing the order of the group of roots of unity in $K$, then
% 	\[
% 		\#\bmat{t\in \cO_{K,S}\xT\mid 1-t\in\cO_{K,S}\xT, t\nin(K\xT)^{2}, K(t^{1/m})\cong L, t\eqv t_{0}\bmod{v}}<+\ift.
% 	\]
% \end{prf}


% \begin{thm}[{\cite[Proposition~5.3, Theorem~5.4]{LV2020}}]
% 	Let $X\ar Y'\ar Y$ be an abelian-by-finite family over a curve $Y$ of genus $g\gqs2$ and $X\ar Y'$ is of relative dimension $d$. If it admits a good model over $\cO_{K,S}$, and $v$ a friendly finite place of $K$, then
% 	\[
% 		\#\bmat{y\in Y(K)\mid \mathrm{size}_{v}(\pi\xI(y))<\frac{1}{d+1}}<+\ift.
% 	\]

% 	Take $X_{q}\ar Y_{q}'\ar Y$ to be the Kodaira-Parshin family, one can conclude $Y(K)$ is finite.
% \end{thm}











% ----------------------------------------------------------------------------------------------------
\bibliographystyle{amsalpha}
\bibliography{../universal-ref}
\end{document}

